\documentclass[border=1pt]{standalone}
\usepackage{mathptmx}  % Times text and math
\usepackage{circuitikz}
\ctikzset{bipoles/length=0.75cm, transistors/scale=0.8,
  inductors/coils=4, current/distance=.5, bipoles/diode/height=.32,
  bipoles/diode/width=.26}
\begin{document}
\fontsize{8}{9.6}\selectfont
% flyback: coordinates are layout units, 1 unit = 1.4 cm across and 1.2 cm up (ieee-single)
\begin{circuitikz}[x=1.4cm, y=1.2cm, line width=0.55pt, cute inductors, european resistors, american voltages]
% ---- wires ----
\draw (0,3) -- (1.3,3);
\draw (1.3,0.6) -- (1.3,0) -- (0,0);
\draw (3,3) -- (4.3,3);
\draw (1.9,1.6) -- (1.9,0) -- (4.3,0);
% ---- Vin (V) ----
\draw (0,3) to[V, l_=$V_\mathrm{in}$] (0,0);
% ---- T (transformer) ----
\draw (1.3,3) to[L] (1.3,1.6);
\draw (1.9,1.6) to[L] (1.9,3);
\draw (1.5643,1.6) -- (1.5643,3) (1.6357,1.6) -- (1.6357,3);
\fill (1.1786,2.5333) circle (1.2pt);
\fill (2.0214,2.0667) circle (1.2pt);
\node at (1.6,3.1917) {$1:n$};
% ---- S (nmos) ----
\node[nmos,bodydiode] (S) at (1.3,1.1) {};
\draw (1.3,1.6) -- (S.D) (S.S) -- (1.3,0.6);
\node[anchor=west] at ([xshift=0.2933cm]S.center) {$S$};
% ---- D (D) ----
\draw (1.9,3) to[D, l^=$D$] (3,3);
% ---- Co (C) ----
\draw (3.4,3) to[C, l_=$C_o$] (3.4,0);
% ---- RL (R) ----
\draw (4.3,3) to[R, l_=$R_L$] (4.3,0);
% ---- annotations ----
\node[currarrow, rotate=0] at (0.65,3) {};
\node[anchor=south, inner sep=0pt] at ([yshift=.1cm]0.65,3) {$i_p$};
\node[currarrow, rotate=0] at (3.2,3) {};
\node[anchor=north, inner sep=0pt] at ([yshift=-.1cm]3.2,3) {$i_D$};
\node at (4.8,1.92) {$+$};
\node at (4.8,1.5) {$v_o$};
\node at (4.8,1.08) {$-$};
% ---- junction dots ----
\foreach \p in {(3.4,0),(3.4,3)} \fill \p circle (1.3pt);
\end{circuitikz}
\end{document}
